% \UseRawInputEncoding
\documentclass[sigplan,nonacm,anonymous]{acmart}
% enable page numbers
% \settopmatter{printfolios=true}
% Optional: Remove the ACM reference between the abstract and the main text.
% Optional: Comment out the CCS concepts and keywords.
\usepackage{algorithm}
\usepackage{algpseudocode}
\usepackage{amsmath}
\usepackage{stfloats}
\usepackage[export]{adjustbox}
\usepackage{booktabs}
\usepackage{tabularx}
\usepackage{graphicx}
\usepackage[justification=centering]{caption}
\usepackage{subcaption}
\usepackage{listings}
\usepackage{url}
\usepackage{hyperref}
\usepackage{xcolor}
\usepackage{circledsteps}
\usepackage[normalem]{ulem} % provides \sout (cross-out)


\DeclareRobustCommand{\legendsquare}[1]{%
  \textcolor{#1}{\rule{1ex}{1ex}}%
}
\usepackage{booktabs, multirow} % for borders and merged ranges
\usepackage{soul}% for underlines
\usepackage{changepage,threeparttable} % for wide tables

% Drop the ACM reference block (explicitly supported by acmart for nonacm papers).
\settopmatter{printacmref=false, printfolios=true}

\newcommand{\column}[2]{%
  \begin{tabular}[b]{@{}c@{}}#1\\#2\end{tabular}%
}
% EK comments (disabled for submission)
\newcommand{\EK}[1]{}
% Auto-generated macros with campaign numbers (capacities, CPU, throughput)
\input{tmp/macros}
\begin{document}
% Flush-bottom default spreads inter-paragraph glue to align column bottoms,
% which pushes figures to the next page and leaves big gaps on earlier pages.
% Ragged-bottom lets columns end naturally, packing more content per page
% without touching any margin / font / spacing parameters.
% \raggedbottom
\title{Astrolabe: Balancing Load in LLM Serving with Context, Knowledge, and Prediction-Guided Scheduling}

\author{Wei Da}
\affiliation{%
  \institution{University of Cambridge}
  \country{United Kingdom}}

\author{Evangelia Kalyvianaki}
\affiliation{%
  \institution{University of Cambridge}
  \country{United Kingdom}}  



% Use the following at camera-ready time to suppress page numbers.
% Comment it out when you first submit the paper for review.
% \thispagestyle{empty}



\begin{abstract}
This paper presents Astrolabe, a \emph{randomized} prediction-guided scheduler for one-shot request dispatch in multi-instance large language model (LLM) serving. Astrolabe improves load balancing without relying on migration-based rebalancing, whose KV-cache transfers can introduce substantial overhead and network contention under high load. Astrolabe combines response-length estimation, per-instance simulation-based latency prediction, and a \emph{power-of-two choices} dispatch policy to improve load balancing and avoid request herding. On the default Llama-2-7B/ShareGPT setup, Astrolabe matches the best load-aware baseline's SLO capacity (31.6 vs 31.5 QPS) while improving latency metrics: 8--36\% lower mean TTFT, 16--77\% lower P99 TTFT, up to 5.6\% lower mean E2E, and $\sim$6$\times$ fewer preemptions once capacity is reached. Under configuration shifts, Astrolabe also improves SLO capacity by up to +6\% on Qwen2-7B and +7.1\% under tight batching, with 6--9\% lower mean E2E under bursty arrivals, while using $\sim$\cpuPoTwoVsFanoutRatio{}$\times$ less per-predictor CPU than full fanout. With migration enabled on A100, Astrolabe outperforms Llumnix by up to 2.6$\times$ in throughput with orders-of-magnitude lower per-token latency at saturation.
\end{abstract}
\keywords{LLM serving, scheduling, load balancing, latency prediction, simulation, KV cache}

\maketitle

\section{Introduction} \label{section:introduction}

The rise of large language models like GPT~\cite{achiam2023gpt}, Llama~\cite{llama}, Gemini~\cite{team2023gemini}, Qwen~\cite{yang2025qwen3}, and DeepSeek~\cite{liu2025deepseek} has revolutionized modern applications, such as chatbots~\cite{achiam2023gpt}, virtual assistants~\cite{guan2023intelligentvirtualassistantsllmbased}, code generation~\cite{jiang2024surveylargelanguagemodels}, and creative writing~\cite{gómezrodríguez2023confederacymodelscomprehensiveevaluation}. This places immense pressure on LLM inference serving systems, which must meet stringent latency requirements for a seamless user experience and maximize throughput to handle growing demand~\cite{shen2025accelgenheterogeneoussloguaranteedhighthroughput, cheng2025scootsloorientedperformancetuning, vllm, orca}. To this end, recent years have seen significant progress in LLM serving systems. Several key techniques, such as continuous batching~\cite{orca}, Paged Attention~\cite{vllm}, chunked prefill~\cite{agrawal2023sarathiefficientllminference}, and FlashAttention~\cite{flashattentionv2}, have been proposed to improve throughput, latency, and utilization.

However, despite the above advances, LLM inference is often characterized as unpredictable~\cite{sun2024llumnix}. For example, the autoregressive nature of LLMs, which generates tokens sequentially based on preceding ones until a stop signal is reached, leads to variable response lengths and decoding steps~\cite{schuurmans2024autoregressivelargelanguagemodels}.
Paged Attention~\cite{vllm}, which dynamically allocates memory resources and allows request preemption, further contributes to the dynamic nature of runtime memory consumption. Decode-step latency also exhibits high variance due to batch-size variation and interference between requests in the same batch~\cite{sun2024llumnix}.

The above uncertainties raise the new challenge of load balancing in multi-instance LLM serving clusters, since widely used runtime metrics such as latency, throughput, and memory no longer accurately represent each request's E2E load. Moreover, scheduled requests with unexpectedly long responses can further increase memory load on overloaded hosts to aggravate load imbalance and block subsequent requests from being launched, known as head-of-line (HOL) blocking~\cite{QLMQueue, ELIS}. In practice multi-instance serving frameworks in production, such as~\cite{nvidiaTritonInference, deepspeed}, employ simple scheduling heuristics like round-robin for request dispatching which offer no guarantee on scheduling performance.  Alternative solutions, such as Llumnix~\cite{sun2024llumnix}, combine dispatching with dynamic rebalancing through live migration and have demonstrated performance improvements. However, migration requires transferring request KV caches~\cite{kang2024gearefficientkvcache} across instances, adding network, memory, and orchestration overhead. Under high load, these costs can accumulate and contend with normal serving traffic, potentially degrading tail latency and throughput.

In parallel to the above optimizations we observe a series of works that aim to mitigate such uncertainties with the help of purpose-built assistant models.
For example, the length of responses and duration of execution can be accurately predicted with a pretrained regression model~\cite{length_estimation}, and via an inference simulation framework~\cite{vidur}.  This presents an unexplored opportunity to improve load balancing in inference systems: if each request's length and latency can be estimated at dispatch time, the scheduler can proactively route long-running requests to less-loaded instances, improving cluster balance before imbalance manifests.

A naive way to exploit such predictions is to query every instance's predictor in parallel and pick the lowest-latency candidate (\emph{full fanout}). This gives the best load-balancing quality but has three practical costs. \emph{First}, each dispatch costs $O(N)$ predictor queries for $N$ instances, which scales poorly in the scheduler control plane. \emph{Second}, simulation runs on CPU to predict the scheduling signal (e.g., latency) for each candidate instance, so aggregate cluster-CPU grows linearly with $N$ per request. \emph{Third}, deterministic fanout creates a \emph{thundering-herd} issue: under short correlated bursts, nearby requests compute similar ``best'' candidates and route to the same instance, overloading it transiently~\cite{mitzenmacher2001,10.1137/S0097539795288490}. These costs matter in any large multi-instance deployment, and become especially acute when scheduler CPU, control-plane bandwidth, or interconnect capacity is limited.

In this paper, we present \emph{Astrolabe}, a randomized prediction-guided scheduler for multi-instance LLM inference clusters. Astrolabe's default dispatch policy is \emph{power-of-two choices} (Po2): for each request, sample two random instances, query their predictors in parallel, and route to the lower-predicted-latency candidate. This yields constant per-request fanout cost, $\sim$\cpuPoTwoVsFanoutRatio{}$\times$ less per-predictor CPU than full fanout in our measurements, and crucially breaks herding by randomizing the probe set so that simultaneous requests are unlikely to sample the same pair of candidates. 

Classical power-of-two choices analysis~\cite{mitzenmacher2001} has been widely applied to load balancing in distributed systems~\cite{sparrow, prequal, hawk, dodoor}, and our measurements confirm the effect in the LLM-serving context. Astrolabe is not tied to a particular cluster topology: its goal is to make high-quality placement decisions before execution, avoiding migration when migration is unavailable, expensive, or harmful due to KV-cache transfer overhead and network contention. Our Llumnix comparison in \S\ref{section:llumnix_comparison} shows that, under high load, migration overhead can collapse serving performance, while Astrolabe maintains higher throughput and lower latency.

Astrolabe's request path is deliberately simple. It first predicts a request's response length from the prompt using a lightweight regression model. The scheduler then samples two candidate instances and queries their Predictor sidecars, each of which simulates the request against its local runtime state and returns an expected latency. The request is dispatched to the candidate with the lower predicted latency. This per-instance view lets Astrolabe improve load balance with one-shot dispatch, avoiding live migration while remaining compatible with existing backend mechanisms. The Predictor interface is pluggable across simulation engines, inference frameworks, and scheduling objectives.

We evaluate Astrolabe against five heuristic/load-aware schedulers and alternative prediction-guided variants on a 12-GPU A30 cluster across ShareGPT, BurstGPT, Llama-2-7B, and Qwen2-7B; we also compare against the official Llumnix release on a 2-node A100 cluster with full migration enabled. On the default A30 Llama-2-7B/ShareGPT setup, Astrolabe preserves Llumnix's SLO capacity (\capPoTwoEst{} vs \capLlumnix{} QPS with no migration), matches full-fanout capacity (\capPoTwoEst{} vs \capFanoutEst{} QPS with no sampling) at 6$\times$ fewer predictor queries, and delivers the lowest TTFT and E2E metrics at 15 of 17 tested QPS points. It also shows up to 97\% lower mean TTFT than round-robin, 8--38\% lower than the best baseline, and $\sim$6$\times$ fewer preemptions at overload ($\sim$300 vs $\sim$1{,}800 at QPS=32). Under configuration shifts, Astrolabe lifts capacity by up to +6\% on Qwen2-7B and +7.1\% on tight batching, tolerates injected prediction noise (at most +6.3\% mean E2E and +3.7\% P99 E2E across the heatmap), and reduces mean E2E by 6--9\% under bursty arrivals. On the two-node A100 cluster, Astrolabe outperforms full Llumnix by up to 2.6$\times$ in throughput and achieves orders-of-magnitude lower per-token latency at saturation. These results show that migration overhead can collapse serving performance under high load, and that prediction-guided one-shot dispatch can provide equal or better load balancing without migration.

This paper makes three main contributions. \emph{First}, we explain why multi-instance LLM serving load is hard to infer from runtime metrics and identify \emph{scheduler herding} issues at \S\ref{section:motivation} and \S\ref{section:scheduling_issues}.
\emph{Second}, we design and implement \emph{Astrolabe}, a one-shot randomized prediction-guided scheduler combining response-length estimation, per-instance simulation-based latency prediction, and a power-of-two choices dispatch policy to achieve high-quality load balancing and request scheduling at \S\ref{section:system_design} and \S\ref{section:implementation}.
\emph{Third}, we evaluate Astrolabe on a 12-GPU A30 cluster and a 2-node A100 cluster across two models, two datasets, and four backend configurations, quantifying prediction accuracy, E2E/TTFT latency, throughput, CPU overhead, robustness under bursty arrivals and injected noise, and a case study comparing prediction-guided dispatch with migration-based rebalancing at \S\ref{section:evaluation}.


\section{Background}\label{section:motivation}
LLMs build on transformer architectures~\cite{transformer}. A prompt is tokenized and processed by transformer blocks to produce intermediate representations (\emph{prefill}); tokens are then generated autoregressively (\emph{decode}) until an end-of-sequence token or max-length limit. Because tokens are reused across steps, attention key/value tensors are cached (\emph{KV cache}~\cite{kang2024gearefficientkvcache}) to avoid recomputation. Prefill is typically compute-bound; decode is memory-bandwidth-bound~\cite{podAttention} and cannot be parallelized across tokens.

Requests are grouped into batches to leverage GPU matrix units; each decoding step generates one token for every in-batch request, and KV cache grows with tokens. Static batching~\cite{lightseqhighperformanceinference, turbotransformersefficientgpuserving} fixes the batch set and under-utilizes slots when shorter requests finish early. Continuous Batching~\cite{orca} addresses this by letting completed requests exit and admitting new successors (FCFS order) up to a maximum batch size. A local scheduler in the inference framework forms these execution batches.

\begin{figure}
    \centering
    \includegraphics[width=0.90\linewidth]{figure/introduction/preemption.png}
    \caption{Paged Attention and Preemption}%
    \Description{Illustration of vLLM's Paged Attention memory blocks and the preemption mechanism when GPU memory becomes insufficient; the newest request is preempted, its KV-cache blocks are released, and it is later recomputed and resumed.}
    \label{fig:pageAttention}%
\end{figure}

A key challenge is memory allocation: response lengths are unknown at runtime, so pre-allocating for max sequence length fragments memory and wastes it on short responses. Paged Attention~\cite{vllm}, implemented in vLLM, partitions GPU memory into fixed-size blocks and tracks each KV cache via a page table, allowing non-contiguous placement (Figure~\ref{fig:pageAttention}). When memory is insufficient for the next decode step, vLLM preempts the newest in-batch request, releases its blocks, and later resumes it by recomputing the KV cache---a key reason for vLLM's practical effectiveness.

Hybrid prefill and decode scheduling is the next challenge. vLLM's original scheduler creates separate prefill-only and decode batches and prioritizes prefill (higher throughput, lower TTFT), but this interrupts decode batches and causes decoding stall bubbles and tail-latency regressions (Figure~\ref{fig:chunked_prefill}).

\begin{figure}
    \centering
    \includegraphics[width=0.90\linewidth]{figure/introduction/chunked_prefill.png}
    \caption{Original vLLM and Chunked Prefill}
    \Description{Comparison of vLLM's original prioritized prefill scheduling with chunked prefill, highlighting how chunked prefill interleaves prefill chunks with decode to reduce decode stalls and tail latency.}
    \label{fig:chunked_prefill}
\end{figure}


Chunked Prefill with a stall-free local scheduler~\cite{agrawal2023sarathiefficientllminference, agrawal2024tamingthroughputlatencytradeoffllm} mitigates this trade-off: prompt processing is split into equal-sized chunks executed across scheduling steps, forming hybrid batches that interleave decode steps with prefill chunks under a token budget, configured as the chunk size. This reduces decode stalls with only minor throughput loss, making it the default in vLLM and SGLang~\cite{sglang}. Prefill-Decode (P-D) disaggregation~\cite{disServe, patel2024splitwiseefficientgenerativellm, qin2024mooncakekvcachecentricdisaggregatedarchitecture} is an alternative that uses separate instances at the cost of cross-instance KV-cache transfers. However, aggregated execution can still be preferable in some settings~\cite{wang2025prefilldecodeaggregationdisaggregationunifying}, and GPU multiplexing~\cite{gao2025duetserveharmonizingprefilldecode} offers similar isolation benefits without disaggregation's KV-cache transfer overhead.

\section{Related Work and Motivation}\label{section:scheduling_issues}

Current LLM inference with dynamic batching and memory allocation, as discussed above, introduces new cloud-based challenges for task scheduling and load balancing. Accordingly, this section first elaborates on these issues and then summarizes applicable techniques, leading to the proposed solution discussed in \S\ref{section:system_design}.

\paragraph{Unpredictability in LLM Serving Scheduling} \label{section:unpredictability} Figure~\ref{fig:llm-inference-workflow} shows a typical architecture for LLM inference serving.  First, a user interacts with an API that serves as the entry point to the inference framework. A local scheduler within the framework then batches these incoming requests for parallel execution. The LLM inference process is composed of a series of layers, primarily involving attention mechanisms and linear projections. To optimize performance, these operations are often accelerated by highly efficient kernel implementations~\cite{flashattentionv2, ye2025flashinfer}. For real-world cloud applications~\cite{wang_gemini_2023,claude, achiam2023gpt}, deploying multiple framework instances is typically necessary to ensure high availability under heavy request volume. This multi-instance setup requires an external global scheduler to distribute incoming requests effectively across the available instances for load balancing.

\begin{figure}
    \centering
    \includegraphics[width=0.90\linewidth]{figure/introduction/inference_workflow.png}
    \caption{LLM Inference Systems}
    \Description{High-level workflow of an LLM inference serving system, including the user/API frontend, local batching scheduler, model execution, and a global scheduler routing requests across multiple instances.}
    \label{fig:llm-inference-workflow}
\end{figure}

LLM inference systems are typically characterized by high unpredictability in both the final outputs and overall system performance over time, as noted in~\cite{sun2024llumnix}. This unpredictability stems primarily from two sources: 1) variable memory demands caused by unknown decoding lengths, which can trigger unexpected preemption; and 2) resource competition, which can create interference between colocated requests. Both effects can degrade runtime performance.

This inherent unpredictability in resource requirements and execution behavior makes dispatching requests across LLM serving instances more challenging than general dynamic task scheduling~\cite{dynamic_task_scheduling}. For example, a standard scheduler like Kubernetes might filter compute nodes based on initially available resources and then select the node with the lowest current utilization. However, such static filtering and simple utilization metrics are often inadequate for LLM inference workloads because request memory demands are dynamic and vary with decoding length. Additionally, the actual performance when executing a request can still be impacted by interference from other colocated requests~\cite{sun2024llumnix}. For example, unexpected preemption due to GPU-memory contention can reset active requests and block subsequent requests. In contrast, traditional cloud-based task execution utilizes isolation mechanisms with static resource allocation (such as Docker container \cite{docker}) that largely reduce the risk of such interference. Furthermore, such isolated and static resource allocation allows task lifetimes to be predicted using historical modeling, as demonstrated by schedulers like Apollo~\cite{Apollo} and LAVA~\cite{ling2025lavalifetimeawarevmallocation}.

These factors make it difficult to accurately predict resource demands and duration of a request on a given instance. Consequently, many existing dispatching strategies for routing requests across LLM instances rely on relatively simple heuristic approaches such as round-robin employed by systems like DeepSpeed-MII~\cite{deepspeed} and Triton Inference Server~\cite{nvidiaTritonInference}. Other frameworks take different approaches. For instance, vLLM often delegates the responsibility of cross-instance routing to the user or a higher-level orchestrator. In contrast, frameworks like SGLang implement their own heuristic-based routers that are tightly coupled with internal optimizations and features, such as RadixAttention. ServerlessLLM~\cite{serverlessLLM} processes LLM requests as serverless functions and schedules them to instances with the least estimated model-loading startup time. LLM-d~\cite{githubGitHubLlmdllmd} serves LLM models in Kubernetes and supports P-D disaggregation with separate customized schedulers. 

Such rule-based heuristic schedulers could be carefully tuned to perform well under specific settings, e.g., configurations, workloads, or models. However, these schedulers are susceptible to performance degradation in dynamic environments and lack the quantifiable metrics needed to clearly explain how the internal parameters/rules impact user-facing metrics. In general, the unpredictability of LLM inference workloads and the lack of accurate performance metrics at dispatch time make it difficult to design effective heuristic schedulers that can consistently achieve good load balancing across instances, especially under high load and with variable workloads. Llumnix~\cite{sun2024llumnix} proposes a dynamic re-balancing approach to mitigate such load imbalance by migrating requests across instances with their KV caches. However, migration can introduce significant overhead and network contention, and can degrade performance and even collapse the entire cluster under high load. This motivates the need for more proactive one-shot scheduling approaches that can improve load balancing without migration. 


% \paragraph{Auto-provisioning in LLM serving} \label{section:scaling} 
% To ensure a stable application and a good user experience in serving clusters with variable workloads, effective resource management must go beyond scheduling to include auto-provisioning to handle peak loads and guarantees performance under established SLOs.
% We notice auto-provisioning and load balancing can be treated as different but complementary problems that are both critical and can be impacted by LLM inference's unpredictability. While load balancing optimizes the spatial placement of requests across existing resources, auto-provisioning manages temporal capacity adjustments. Although highly desirable and well-studied in other cloud applications, auto-provisioning is often overlooked in LLM serving systems. Asynchronous cold starts in serverless computing \cite{serverless_real_world} exemplify the issue. While a new instance is being provisioned and initialized, incoming requests continue to be routed to existing instances, exacerbating tail latency. The effect is even more pronounced for LLM serving. In addition to new arrivals, ongoing inference requests also continue generating tokens and consuming additional memory on already overloaded instances. This prolongs latency and increases load, even as the newly provisioned instances remain underutilized. Finally, the resulting load imbalance wastes resources and degrades serving performance.

% Llumnix~\cite{sun2024llumnix} can mitigate such imbalances through dynamic re-balancing. Initially, it dispatches requests using a heuristic scheduler. Subsequently, Llumnix performs continuous dynamic load rebalancing across instances by migrating active requests along with their KV caches. While dynamic rebalancing can mitigate runtime load imbalances, it necessitates transferring associated token caches, which consumes inter-GPU network bandwidth and introduces transient memory duplication that can impact model parallelism performance. This resource contention can even intensify in either: 1) when auto-provisioning is triggered due to cold start issues, to re-balance load between heavily utilized existing hosts and newly provisioned ones; or 2) when other network-heavy features are applied, such as P-D disaggregation and tensor/pipeline parallelism~\cite{shoeybi2020megatronlmtrainingmultibillionparameter, huang2019gpipeefficienttraininggiant}. 

% Despite the significant challenges caused by unpredictability in LLM scheduling and provisioning, recent research offers promising mitigation techniques that can improve scheduler design, which we explore below.


\paragraph{Offline Performance Simulator} \label{section:performance simulation}

With the rapid development and expansion of new LLM models, devices, and use cases, setting up an appropriate LLM serving cluster has become a critical yet challenging task for developers, particularly due to the high trial-and-error costs associated with GPU pricing. To address this, Vidur, an LLM serving cluster simulation framework, has been proposed~\cite{vidur}. It aims to reduce potential hardware costs when searching for the optimal cluster configuration based on a given model, trace, and service-level objective (SLO) requirements. The insight behind Vidur is that since the local scheduler and batching logic are deterministic, and response lengths are typically available in the replay trace, it is feasible to simulate the entire replay process if the execution time for each batch can be inferred. Additionally, by profiling low-level operators like attention and linear projection for a specific GPU, Vidur trains linear models to interpolate execution times for various batches. It reports less than 9\% error for key metrics such as throughput and latencies. Other frameworks, such as LLMServingSim~\cite{cho2025llmservingsim}, Frontier~\cite{feng2025frontier}, Revati~\cite{agrawal2026revati}, and LLM-Emu~\cite{da2026llmemunativeruntimeemulation} have also been proposed to simulate LLM serving clusters with different areas of focus.

\paragraph{Response Length Prediction} \label{section:length_prediction} Prior work predicts unknown response lengths to help \emph{local} schedulers: Sequence Scheduling~\cite{length_estimation} fine-tunes another LLM to predict response length from the input instruction and batches similar-length requests to achieve +85\% throughput. LightLLM~\cite{lightLLM} uses historical distributions to estimate the expected length to avoid preemption. The ELIS~\cite{ELIS} local scheduler uses a regression model to estimate the length to construct batches with shorter completion times.

Length estimation has also been utilized for global scheduling, primarily as a filtering mechanism prior to scheduling tasks. For instance, TetriServe~\cite{TetriServe} deploys a lightweight length-estimation model to filter out decoding candidate instances without sufficient GPU memory to accommodate the estimated tokens and then applies the power-of-two method~\cite{poweroftwo} to select decoding instances with fewer pending requests. DynamoLLM~\cite{dynamollm} classifies requests into three distinct pools based on their estimated lengths, scheduling them separately with associated host pools. To summarize, our insight is that by combining simulation techniques with length prediction, the unpredictability can be substantially reduced, which in turn creates an opportunity to improve multi-instance LLM serving performance.

% Figure 4 declared before §4 so the float can land at bottom of p4's right
% column instead of being deferred to top of p5 and leaving p4 short.
\begin{figure}[!tb]
    \centering
    \includegraphics[width=0.90\linewidth]{figure/block_diargram.png}
    \caption{Astrolabe Architecture}
    \Description{System architecture of Astrolabe, showing four components (query length tagger, global scheduler, predictor sidecars attached to each model instance, and the inference framework backend) and the flow from incoming requests through prediction, scheduling, execution, and returning results.}
    \label{fig:scheduling_framework}
\end{figure}

\section{Astrolabe System Design} \label{section:system_design}

In this section, we present \emph{Astrolabe}, a prediction-guided task scheduler that leverages query context and cluster knowledge for load balancing in LLM serving clusters. Astrolabe builds on the key insight from \S\ref{section:motivation} that LLM serving scheduling is challenging due to its uncertainty and couples runtime prediction with a randomized dispatch policy to achieve high-quality load balancing.

Astrolabe, as shown in Figure~\ref{fig:scheduling_framework}, comprises four services: a query length tagger, a global scheduler, per-instance Predictor sidecars, and inference framework backends. The query length tagger is the entry point and predicts the anticipated response length for each request (\Circled{1}). Requests are then forwarded to the global scheduler (\Circled{2}), which selects a target model instance.
For each incoming request, the global scheduler issues \texttt{predict} RPCs to a sampled subset of Predictor sidecars (\Circled{3}) (power-of-$k$ choices; $k{=}2$ by default, and $k{=}N$ recovers full fanout). Each Predictor queries its colocated backend via a \texttt{status} API to obtain the instance's current runtime state (e.g., running/waiting requests, batching state, and memory availability), and then runs a lightweight simulation to estimate the request's target metric (e.g., E2E latency) (\Circled{4}). The global scheduler aggregates these per-instance predictions and dispatches the request to the selected backend instance according to the configured scheduling policy (\Circled{5}). Finally, the inference backend executes the request on the GPU and returns the response. Astrolabe is designed as an add-on scheduling layer and is agnostic to models, hardware, inference frameworks, and dispatch policies.


\subsection{Model Instance} \label{section:model_instance}

In Astrolabe, we define a model instance as the collection of services deployed on a GPU host that uses one or more GPUs for model execution. It consists of two primary components: the main services of the inference framework and a sidecar service called the Predictor. Astrolabe is designed to be framework-agnostic and can support different frameworks with dynamic batching, such as~\cite{vllm, orca, deepspeed}, which are already supported by Vidur.
Each inference framework used by Astrolabe can be integrated by supporting a new \texttt{status} API to export its internal status, such as request lists and GPU memory blocks, for metric prediction.

The Predictor's main role is to predict key performance metrics, such as E2E latency or TTFT, for incoming requests.
As shown in Figure~\ref{fig:scheduling_framework}, the Predictor service runs locally on each instance, aggregates runtime data exported by the \texttt{status} API, and transforms it into a scalar metric prediction via the \texttt{predict} API, which is then consumed by the global scheduler.
In cases where a request's decoded length exceeds its predicted length at runtime, the simulator extends the current prediction by assuming the request will complete within the next 10 tokens by default.

Astrolabe is also designed to work with different simulation approaches. In our current prototype implementation of the Predictor's simulator, we integrate the Vidur~\cite{vidur} simulation framework in the following way. We redesign the Vidur simulator for single-instance prediction and encapsulate it within the Predictor service API. As illustrated in Figure~\ref{fig:scheduling_framework}, this simulation involves a two-stage process. First, a local scheduler simulator models the batching strategy for a given inference backend. Second, a linear model predicts the execution times for the batches generated in the first stage. Since the model is static and inputs are received on demand from the \texttt{status} API, the Predictor service does not maintain scheduling state and can be replicated per instance to reduce simulation-related overhead through parallelization. Furthermore, the Predictor service is designed with an extensible interface, which could be implemented by alternative simulation frameworks while keeping core strategies.

\subsection{Global Scheduler} \label{section:schedulers}

\begin{algorithm}[t]
\caption{Astrolabe one-shot dispatch (default: latency)}
\label{alg:astrolabe_dispatch}
\begin{algorithmic}[1]
\Require Request $r$ with prompt $p$; instances $\mathcal{I}=\{i_1,\dots,i_N\}$; sample size $k$ (default $2$); timeout $\tau$
\Ensure Dispatched instance $i^\star$
\State $\hat{L} \gets \textsc{LengthTagger}(p)$
\State $\mathcal{I}' \gets \textsc{Sample}(\mathcal{I}, k)$
\State $\mathsf{cand} \gets \emptyset$
\ForAll{$i \in \mathcal{I}'$ \textbf{in parallel}}
    \State $s_i \gets \textsc{GetStatus}(i)$
    \State $\hat{\ell}_i \gets \textsc{PredictMetric}(i, r, \hat{L}, s_i)$
    \If{$\hat{\ell}_i$ returned before $\tau$}
        \State $\mathsf{cand} \gets \mathsf{cand} \cup \{(i,\hat{\ell}_i)\}$
    \EndIf
\EndFor
\If{$\mathsf{cand}=\emptyset$}
    \State \textbf{reject} $r$ and \textbf{throw exception}
\EndIf
\State $i^\star \gets \arg\min_{(i,\hat{\ell}_i)\in \mathsf{cand}} \hat{\ell}_i$
\State \textsc{Dispatch}$(i^\star, r)$
\end{algorithmic}
\end{algorithm}


The global scheduler service is stateless and lightweight without maintaining a global table of instance statuses. Instead, as illustrated in Figure~\ref{fig:scheduling_framework}, upon each request arrival, the scheduler issues a \texttt{predict} RPC to sampled Predictor sidecars and then dispatches the request to the instance with the lowest predicted metric. In this paper, we use E2E latency as the default target metric, while keeping the interface configurable to support other objectives if needed.

This differs from centralized predictors that need a global per-instance status view. In Astrolabe, prediction runs on each instance in parallel using its local state, and each Predictor returns only a scalar metric (e.g., predicted latency). Decision time is bounded by the slowest Predictor reply, while avoiding the cost of exporting heavy instance state (e.g., request lists) to a central simulator. Because each Predictor simulates the request on top of its execution queue, the prediction reflects both running and pending load.

To avoid the $O(N)$ per-dispatch messaging cost of full fanout and the scheduler-herding issue (\S\ref{section:introduction}), Astrolabe's default samples $k{=}2$ random instances (power-of-two choices) and queries only those Predictors. Setting $k{=}N$ recovers full fanout, and we sweep $k$ in \S\ref{section:po2_ablation} to study the trade-offs. The dispatch procedure is shown in Algorithm~\ref{alg:astrolabe_dispatch}.

Furthermore, Astrolabe does not require instance-side features such as live migration for dynamic rebalancing~\cite{sun2024llumnix}. Such features can be complementary in some deployments, but they introduce additional orchestration and state-transfer overhead. In this work, we focus on prediction-guided one-shot dispatch and show that it can match or outperform migration-based rebalancing.

\subsection{Query Length Tagger}\label{section:query_tagger}
The query length tagger is an online service that estimates the response length of each request from its prompt and the target serving model. It runs in parallel with request handling and uses a lightweight proxy model to keep overhead low; the estimator is modular and can be replaced with alternatives such as the model-free, sampling-based approach used in LightLLM~\cite{lightLLM}. The predicted length is used only for metric prediction and scheduling decisions and does not change the inference parameters and results.

\section{Implementation} \label{section:implementation}

Astrolabe implements the global scheduler, per-instance Predictor sidecars, and the length-tagger service (\S\ref{section:system_design}) as lightweight FastAPI~\cite{tiangoloServerWorkers} microservices with a simple RPC/REST interface that integrates with common LLM serving stacks (e.g., vLLM~\cite{vllm}, SGLang~\cite{sglang}). The prototype integrates with vLLM 0.7.2 and comprises about 4{,}000 LoC, released at \textcolor{blue}{\url{https://anonymous.4open.science/r/Astrolabe-BA60/README.md}}.

\paragraph{Simulation-based prediction}
The Predictor is the per-instance sidecar described in \S\ref{section:model_instance}; internally it invokes a backend-specific simulator. Our prototype reuses Vidur's offline simulator codebase~\cite{vidur} as the starting point, due to its high-fidelity modeling, broad backend support, and extensibility (Sarathi-Serve, vLLM, LightLLM), and substantially adapts it for online, per-request use. Vidur's offline pipeline has inefficiencies that are negligible offline but become critical for real-time prediction (e.g., object duplication, suboptimal \texttt{list.pop(0)} calls), so we rewrote the primary simulation hot paths and added a per-Predictor cache that memoizes latency predictions per (batch size, token count), substantially reducing the per-request simulation cost.

\paragraph{Framework Integration and Global Scheduler Implementation} \label{section:baselinescheduler}
Astrolabe is decoupled from any specific inference framework (\S\ref{section:model_instance}): integrating a new backend requires only exposing the read-only \texttt{status} API and implementing a local-scheduler simulator. Extending to new serving features (e.g., P-D disaggregation) similarly requires updating only the simulator and, if needed, the state exported via \texttt{status}.

Astrolabe's global scheduler exposes a configurable target metric and dispatch policy, supporting single- and multi-SLO strategies~\cite{chen2025slosserveoptimizedservingmultislo}. Our prototype selects the instance with the lowest predicted E2E latency. For evaluation we implement five baselines within the same framework:
\begin{itemize}
    \item \emph{Random}: selects an instance uniformly at random.
    \item \emph{Round-Robin}: cyclic dispatch, used by production-grade serving systems~\cite{nvidiaTritonInference, rayServeScalable}.
    \item \emph{Min QPM (Queries Per Minute)}: LiteLLM's~\cite{litellm} default policy to dispatch based on QPM as a proxy for load.
    \item \emph{INFaaS++}: the Llumnix-optimized variant of INFaaS~\cite{sun2024llumnix,infaas}, with load as $usedMemory/batchSize$.
    \item \emph{Llumnix-}: Llumnix's~\cite{sun2024llumnix} heuristic dispatcher extending INFaaS++ with prefill memory correction as:
    
    $load = (usedMemory + prefillMemory) / batchSize$.
    Live migration is disabled in the main evaluation and full migration-enabled Llumnix is evaluated in~\S\ref{section:llumnix_comparison}.
\end{itemize}

\paragraph{Length Estimation Model}
For the query length tagger (\S\ref{section:query_tagger}), we fine-tuned a lightweight RoBERTa-base~\cite{roberta} regression model (125M parameters); similar BERT-based estimators are widely adopted in serving systems for auxiliary tasks ranging from model routing~\cite{wang2025reasonsemanticroutervllm, ong2025routellmlearningroutellms, ding2025bestrouteadaptivellmrouting} to output-length prediction~\cite{ELIS, s3}.~\footnote{We initially considered the 7B length model from Sequence Scheduling~\cite{length_estimation}, but its training duration and serving overhead were prohibitive.}

\section{Evaluation} \label{section:evaluation} 
We evaluate Astrolabe against five baseline schedulers (\S\ref{section:baselinescheduler}). We introduce the experimental setup at \S\ref{section:experiment_setting}, validate prediction accuracy at \S\ref{section:predicting accuracy}, and quantify E2E gains in \S\ref{section:request_performance} (including GPU memory dynamics and a head-to-head comparison against Llumnix at \S\ref{section:llumnix_comparison}). We then present additional analyses on power-of-$k$ choices, with $k{=}2$ by default (\S\ref{section:po2_ablation}), oracle-length upper bounds (\S\ref{section:oracle_bounds}), burstiness (\S\ref{section:burstiness}), prediction-noise sensitivity (\S\ref{section:error_sensitivity}), and CPU overhead (\S\ref{section:cpu_overhead}).


\begin{figure*}[!tb]
    \centering
    \includegraphics[width=0.87\linewidth]{figure/exp_plots/cluster_metrics/profiling.png}
    \caption{Latency Prediction Metrics}
    \Description{Metrics summarizing latency prediction accuracy and scheduling effectiveness, including prediction error and the rank of the instance selected by the scheduler among all candidates.}
    \label{fig:latency_prediction}
\end{figure*}

\subsection{Experimental Setup} \label{section:experiment_setting}

\paragraph{Testbed} We evaluate Astrolabe on CloudLab~\cite{cloudlab} using two clusters. The primary cluster consists of 12 d7525 nodes, each equipped with two 16-core AMD 7302 CPUs (3.00\,GHz), 128\,GB ECC memory, one NVIDIA A30 GPU (24\,GB), and 25\,Gbps effective bandwidth due to mandatory best-effort network provisioning in CloudLab. This single-GPU-per-host configuration is used for the majority of experiments with implementations and baselines at (\S\ref{section:implementation}).

Additionally, we deploy a 2-node CloudLab d8545 cluster with 4$\times$ A100-40GB GPUs per node and 100\,Gbps NICs. This 4-GPU-per-node configuration is comparable to Llumnix's original testbed~\cite{sun2024llumnix}. The multi-GPU cluster enables: (1) comparison with full Llumnix migration (\S\ref{section:llumnix_comparison}) with KV-cache transfers; and (2) evaluation of prediction accuracy on larger models with tensor parallelism (\S\ref{section:predicting accuracy}). We run the length-estimation model on a dedicated host with an L40 GPU to process datasets offline. Such lightweight BERT-based models are capable of online execution, as evidenced by their adoption in recent production LLM services~\cite{wang2025reasonsemanticroutervllm, ong2025routellmlearningroutellms}, and we move them offline to reduce interference with the main evaluation and to enable faster iteration during development.

\paragraph{Data, Configuration and Model}
We use ShareGPT~\cite{shareGPTData}, a real-world dataset with 50K conversations for the main experiments. During experiments, prompts are sent to the global scheduler according to a Poisson arrival process at varying rates, referred to below as external QPS.

When serving LLMs on vLLM, we use FlashInfer~\cite{ye2025flashinfer} as the attention library and greedy decoding~\cite{jm3} with temperature 0. The maximum request batch size is 48, and the Chunked-Prefill chunk size is 512; this combination gives the best performance among tested configurations. To reduce simulation overhead, each instance's node runs 16 Predictor replicas to simulate and predict latencies in parallel, which reduces total scheduling overhead by up to 50\%.

Since each A30 test node has a single GPU and limited cross-node bandwidth, co-serving across GPUs using tensor or pipeline parallelism~\cite{narayanan2021efficientlargescalelanguagemodel} is infeasible. We therefore use a medium-sized model that fits within a single node: Llama-2-7B~\cite{llama} with FP16 precision. We choose Llama-2 to stay aligned with Vidur~\cite{vidur}, which lets us reuse Vidur's published model configuration when profiling operators for our simulator; the Llama family is also widely used in prior AI/systems papers. The model weights occupy 12.5\,GB of GPU memory, split into 1056 vLLM memory blocks for KV cache. For tensor-parallel prediction accuracy, we also evaluate Llama-2-70B on the A100 cluster (\S\ref{section:predicting accuracy}) and observe comparable latency-prediction accuracy (MAPE 16.6\%), suggesting that Astrolabe's predictor generalizes to production-scale deployments. Vidur~\cite{vidur} also reports that simulation accuracy can improve with larger models because framework CPU overhead accounts for a smaller fraction of E2E latency.

\subsection{Prediction Accuracy} \label{section:predicting accuracy}

\begin{table}[tb]
    \caption{Response Length Prediction}
    \centering
    \begin{tabular}{ccc}
    \toprule
    Metric & RoBERTa Regressor & Prompt-based LLM  \\ \midrule
    MAE & 78.755 & 62 \\
    MAPE & 24.4\% & Not reported \\
    Acc-50 & 69.93\% & 59\% \\
    Acc-100 & 77.15\% & 81\% \\ 
    \bottomrule
    \end{tabular}
    \caption*{Acc-X: percentage of data points with absolute error of predicted tokens less than X.}
    \label{table:length_prediction}
\end{table}%

\paragraph{Length Prediction.}
We divide the first 50k ShareGPT records into 40k training and 10k evaluation samples. The labels are actual output lengths generated by the served model, and we use them to fine-tune a RoBERTa-base regression model. The resulting 125M-parameter model reaches MAE 78.8 and MAPE 24.4\% on the 10k test split, which is comparable to the 7B prompt-based model reported in Sequence Scheduling~\cite{length_estimation} (Table~\ref{table:length_prediction}; updated numbers from~\cite{githubGitHubZhengzangwSequenceScheduling}). Offline evaluation on 10k requests completes in 4.8\,s in PyTorch, consistent with online overheads reported for comparable BERT-based estimators in other runtime systems ($S^3$: 3.7\,ms~\cite{s3}; ELIS: 11.4\,ms incl.\ batching~\cite{ELIS}).


\paragraph{Simulation-based Metric Prediction.}
Figure~\ref{fig:latency_prediction} assesses the runtime simulator's accuracy and its effectiveness for scheduling. Astrolabe processes requests with real lengths under QPS=20--36 using a random scheduler, recording predicted and actual latencies. Online requests are sampled at 1\%. For each sampled request, we broadcast the request to all instances to collect predicted and actual latencies, then mark the minimal-predicted-latency instance as selected.


% \begin{figure}
%     \includegraphics[width=\linewidth]{figure/exp_plots/cluster_metrics/cdf.png}
%     \caption{CDF for Latency Metrics}
%     \label{fig:cdf}
% \end{figure}

The top row in the figure reports the prediction error rate. We observe that enabling Chunked Prefill yields significantly lower error than vLLM's default Prioritized Prefill. This is because Prioritized Prefill causes decoding stalls triggered by future request arrivals, and such events are unknown at the moment of prediction. Consequently, the simulator cannot model these random interruptions. Chunked Prefill mitigates this by interleaving prefill and decode phases as explained in \S\ref{section:motivation}, aligning runtime behavior with the simulator's linear model. The middle and bottom rows further validate the scheduling effectiveness. The scatter plots in the middle row show that the latency of the selected instance (red line) consistently tracks the lower envelope of the latency distribution (blue dots) and remains significantly below the cluster mean (black line). This demonstrates that even when the scheduler does not select the absolute optimal instance, it consistently avoids overloaded nodes and routes to a high-performance candidate. The bottom row displays the 'Real Latency Position', denoting the rank of the selected instance among all candidates (where 1 is optimal). While the absolute prediction error ranges from 10 to 15\%, Astrolabe mainly relies on \emph{relative} latency ranking for instance selection: the scheduler maintains a high probability (40 to 80\%) of selecting the Rank-1 instance. This is because a large fraction of prediction error manifests as near-constant system noise (e.g., RPC overhead) that affects all instances similarly and therefore does not change their relative ordering.

To validate that Astrolabe's simulation accuracy extends to larger models with tensor parallelism, we evaluate Llama-2-70B on the 2-node A100 cluster with TP=4 (2 instances, 4 GPUs each). At QPS=3 with a 10\% sampling rate, the simulator achieves a latency-prediction error of 16.6\%. This 16.6\% is well within Astrolabe's robustness envelope: \S\ref{section:error_sensitivity} shows the scheduler tolerates up to $\pm$100\% independent Gaussian noise on simulator latency with at most +6.3\% mean E2E degradation and +3.7\% P99 E2E degradation across the heatmap, because Po2 dispatches on \emph{relative} ranking and common-mode error cancels as discussed later. Vidur~\cite{vidur} additionally reports that simulator accuracy may improve at larger model scales because unmodeled per-request CPU overhead becomes a smaller fraction of E2E latency.

\begin{figure*}[!tb]
        \centering
          \includegraphics[width=1.0\linewidth]{figure/exp_plots/cluster_metrics/qps.png}
        \caption{Request Metrics Under Different QPS}
        \Description{Key serving metrics (e.g., E2E latency, TTFT, tail latency, throughput, and scheduling overhead) measured across different external query-per-second (QPS) levels for Astrolabe and baseline schedulers.}
        \label{fig:request_metrics}
\end{figure*}

\subsection{Request Latency Performance} \label{section:request_performance}
Figure~\ref{fig:request_metrics} evaluates Astrolabe ($k{=}2$, RoBERTa lengths) against five non-predictive baselines as discussed in \S\ref{section:baselinescheduler} across an integer-QPS sweep. We report E2E latency, TTFT, scheduling overhead, serving \emph{capacity} (max QPS meeting SLO P99 TTFT\,$\le$\,\sloSec\,s, 0.1-QPS resolution), and request throughput. Llumnix- is the primary non-predictive comparator. Table~\ref{tab:request_metrics_snapshot} (snapshot at QPS=20, 24, 28, 32, 36) quantifies the latency metrics in more detail and shows Astrolabe's gain over the best non-Astrolabe scheduler at each QPS.

\begin{table*}[!tb]
\caption{Per-scheduler TTFT(ms)/E2E(s); last row is Astrolabe gain vs best non-Astrolabe (negative=worse).}
\centering
\begingroup
\setlength{\tabcolsep}{3.6pt}
\renewcommand{\arraystretch}{1.10}
\resizebox{\textwidth}{!}{%
\begin{tabular}{l ccccc c ccccc c ccccc c ccccc}
\toprule
& \multicolumn{5}{c}{TTFT mean (ms)} & & \multicolumn{5}{c}{TTFT P99 (ms)} & & \multicolumn{5}{c}{E2E mean (s)} & & \multicolumn{5}{c}{E2E P99 (s)} \\
\cmidrule(lr){2-6} \cmidrule(lr){8-12} \cmidrule(lr){14-18} \cmidrule(lr){20-24}
QPS $\to$ & 20 & 24 & 28 & 32 & 36 & & 20 & 24 & 28 & 32 & 36 & & 20 & 24 & 28 & 32 & 36 & & 20 & 24 & 28 & 32 & 36 \\
\midrule
Astrolabe & \textbf{67} & \textbf{71} & \textbf{88} & \textbf{2586} & \textbf{18107} & & \textbf{153} & \textbf{161} & \textbf{355} & \textbf{12152} & \textbf{45506} & & \textbf{8.8} & \textbf{9.4} & \textbf{10.3} & \textbf{14.7} & \textbf{30.6} & & 27.3 & \textbf{29.2} & 32.0 & \textbf{41.6} & \textbf{67.9} \\
Llumnix-  & 96 & 111 & 133 & 2935 & 19707 & & 605 & 684 & 759 & 14505 & 54989 & & 8.8 & 9.5 & 10.4 & 15.3 & 32.4 & & \textbf{27.2} & \textbf{29.2} & 31.8 & 43.3 & 75.8 \\
INFaaS++  & 106 & 131 & 149 & 4111 & 22456 & & 630 & 745 & 840 & 43176 & 182614 & & 8.8 & 9.5 & 10.4 & 15.8 & 34.7 & & \textbf{27.2} & \textbf{29.2} & \textbf{31.6} & 58.6 & 193.9 \\
RR        & 99 & 279 & 1381 & 8113 & 20857 & & 1320 & 5719 & 12735 & 36266 & 64683 & & 9.0 & 9.9 & 12.1 & 20.0 & 33.4 & & 28.5 & 31.4 & 38.3 & 56.6 & 83.1 \\
MinQPM    & 127 & 624 & 2153 & 8538 & 21559 & & 2154 & 9118 & 17788 & 36160 & 71727 & & 9.1 & 10.4 & 12.9 & 20.6 & 34.0 & & 28.9 & 33.9 & 41.0 & 56.6 & 86.6 \\
Random    & 158 & 501 & 2279 & 9309 & 21969 & & 2529 & 9485 & 19881 & 53212 & 81112 & & 9.1 & 10.3 & 13.2 & 21.2 & 34.4 & & 29.1 & 33.2 & 42.8 & 67.8 & 96.6 \\
\midrule
Astrolabe gain & $+30.2\%$ & \textbf{+36.0\%} & $+33.8\%$ & $+11.9\%$ & $+8.1\%$ & & $+74.7\%$ & \textbf{+76.5\%} & $+53.2\%$ & $+16.2\%$ & $+17.2\%$ & & $+0.0\%$ & $+1.1\%$ & $+1.0\%$ & $+3.9\%$ & \textbf{+5.6\%} & & $-0.4\%$ & $+0.0\%$ & $-1.3\%$ & $+3.9\%$ & \textbf{+10.4\%} \\
\bottomrule
\end{tabular}%
}
\endgroup
\label{tab:request_metrics_snapshot}
\end{table*}

INFaaS++ performs well at low and moderate load, but collapses with increasing QPS because its memory-load heuristic relies entirely on the current memory state and cannot anticipate the impact of incoming requests. Llumnix- performs better than INFaaS++ at high load by adding a correction term for known prefill demand from requests, but it still misses the decoding dynamics of requests. Astrolabe's predictor simulates the full request lifecycle and therefore can route around overloaded nodes before they collapse, yielding significant latency reductions and throughput improvements without sacrificing capacity.

Across the 17-point sweep, Astrolabe's peak reductions over the best load-aware baseline are 38\% mean TTFT, 78\% TTFT-P99, 7\% mean E2E, and 12\% E2E-P99, with up to 5\% higher throughput; vs the worst baseline they reach 97\% TTFT reduction and 53\% higher throughput. Astrolabe's per-request overhead stays under 3\% of E2E within capacity; per-predictor CPU is $\sim$\cpuPoTwoVsFanoutRatio{}$\times$ lower than full fanout (\S\ref{section:cpu_overhead}), which probes all instances before dispatching. Astrolabe's tail-TTFT gain over other schedulers follows a U-shaped curve across QPS: it shows strong gains at both low load and overload, but falls behind Llumnix- from QPS 29 to 31. This is because at low load, the cluster is mostly idle and the simulator can accurately predict the latency of incoming requests, so Astrolabe can route around even transient stragglers; at high load, the herding effect of overloaded nodes becomes severe and is mitigated by Astrolabe's randomized dispatching and accurate anticipation of request dynamics to avoid collapse and excessive preemption.
At moderate load, the misprediction of request dynamics causes some suboptimal routing decisions that specifically affect tail latency, which makes the capacity gain less pronounced as judged by the TTFT-P99 metric (31.6 vs 31.5). However, the mean TTFT and E2E metrics are less sensitive to such tail events, so they still show consistent gains across the sweep. This suggests that even when the simulator's accuracy is not perfect, Astrolabe's overall scheduling decisions are still improved compared to non-predictive baselines, especially in terms of mean latency and throughput.

Figure~\ref{fig:resource} also reports cross-instance variance of free GPU blocks and cumulative preemptions (both are Gaussian-smoothed with the same parameters to improve readability). Below capacity (QPS\,$\le$\,29) Astrolabe keeps variance low and preemptions near zero. The gap explodes at capacity: at QPS=32 Llumnix- accumulates $\sim$1.8k preemptions vs Astrolabe's $\sim$300 ($\sim$6$\times$ fewer); at overload (QPS 33--36) Llumnix- saturates around $\sim$2.1k preemptions while Astrolabe's count grows near-linearly and stays under 350. This aligns with the sustained TTFT/E2E lead past capacity.

\begin{figure*}[!tb]
        \centering
        \includegraphics[width=0.90\linewidth]{figure/exp_plots/cluster_metrics/linear.png}
         \caption{GPU memory blocks and preemption dynamics across QPS.}
         \Description{These time-related figures have been smoothed using a unified Gaussian filter to reduce noise and enhance readability.}
         \label{fig:resource}
\end{figure*}

% End-to-end comparison with Llumnix on A100 (Astrolabe as default)
\subsection{End-to-end Comparison with Llumnix}\label{section:llumnix_comparison}

We compare Astrolabe with full Llumnix~\cite{sun2024llumnix}, which performs dynamic load rebalancing via live KV-cache migration. Our A30 cluster (single-GPU-per-host, 25\,Gbps Ethernet) is unsuitable for migration due to network constraints. We therefore use a 2-node cluster with 4$\times$A100-40GB GPUs per node, analogous to Llumnix's original 4-node A10 testbed. Llumnix moves KV-cache via Ray RPC by default and does not support chunked prefill (CP) in the evaluated configuration. To keep the comparison faithful to the released system, we keep Llumnix's default Ray-RPC migration path and do not retune its migration policy for our workload; we change only deployment parameters needed to run the same model and instance count on the A100 cluster. To isolate the scheduler's contribution from chunked-prefill backend effects, we evaluate Astrolabe with and without CP against Llumnix across an identical QPS sweep. With larger A100 memory (40\,GB vs 24\,GB), we raise the batch size to 96. All configurations serve Llama-2-7B across 8 instances (4 per node, TP=1 for each), using FlashInfer as backend.

\begin{figure}[!tb]
    \centering
    \includegraphics[width=0.90\linewidth]{figure/revision/llumnix_comparison_v2_po2.png}
    \caption{Astrolabe with and without chunked prefill vs full Llumnix on the A100 cluster}
    \Description{Throughput, mean request latency, mean token latency, and scheduling overhead versus QPS for Astrolabe with and without chunked prefill and full Llumnix on the A100 cluster.}
    \label{fig:llumnix_comparison}
\end{figure}

Figure~\ref{fig:llumnix_comparison} reports throughput, mean of request latency, token latency, and scheduling overhead. Astrolabe incurs a higher mean scheduling overhead than Llumnix at low load (23--29\,ms vs 5--7\,ms at QPS\,$\le$\,19) because it runs predictor queries and simulation before dispatch, but the overhead stays in the tens of milliseconds. Through QPS=19, both systems serve comparably; CP gives Astrolabe a small edge (5.5\,s vs 6.3\,s mean E2E at QPS=19), while the CP and no-CP Astrolabe curves remain within $\sim$10\% of each other. The high-load gap is therefore not explained by CP alone: Astrolabe without CP already outperforms Llumnix substantially. Llumnix collapses at QPS=20, where mean request latency jumps from 6.3\,s to 106\,s and mean token latency from 21\,ms to 1.15\,s. This is consistent with migration-induced PCIe/NIC contention: repeated KV-cache transfers can contend with normal serving traffic, allowing migration overhead to accumulate and dominate E2E latency. At QPS=36, the direct no-CP comparison already favors Astrolabe: Astrolabe without CP delivers 2.57$\times$ the throughput (18.7k vs 7.3k\,tok/s), 36$\times$ lower mean request latency (8.2\,s vs 300\,s), and 123$\times$ lower mean token latency (26\,ms vs 3.2\,s) than Llumnix. Enabling CP widens these to 2.58$\times$/40$\times$/135$\times$.

% k-tunable ablation — power-of-choices family (Po2 default, Po4, Po8, Fanout)
\subsection{\texorpdfstring{$k$-tunable Dispatch: from Po2 to Fanout}{k-tunable Dispatch: from Po2 to Fanout}}\label{section:po2_ablation}

Astrolabe's core dispatch policy is \emph{power-of-$k$ choices}: for each request, sample $k$ random instances, query each Predictor in parallel, and route to the one with the lowest predicted latency. Increasing $k$ trades higher control-plane fanout cost (linear in $k$) for better load-balancing quality and stronger mitigation of herding~\cite{mitzenmacher2001,10.1137/S0097539795288490}. We evaluate $k \in \{2, 4, 8, 12\}$; $k{=}12$ recovers full Fanout.

\begin{figure}[!tb]
    \centering
    \includegraphics[width=0.90\linewidth]{figure/revision/po2_comparison.png}
    \caption{Power-of-$k$ choices ablation}
    \Description{Capacity, E2E latency, throughput, and scheduling overhead for power-of-k choices variants from k=2 to full fanout.}
    \label{fig:po2_comparison}
\end{figure}

Figure~\ref{fig:po2_comparison} shows capacity, mean/P99 E2E latency at each scheduler's own capacity, throughput, and per-request scheduling overhead. Capacity is effectively tied from $k{=}2$ upward: Po2-est reaches \capPoTwoEst{} QPS vs Po4-est \capPoFourEst{} (the maximum), Po8-est \capPoEightEst{}, and Fanout-est \capFanoutEst{}, all within a $\pm$0.3\,QPS band at SLO\,=\,\sloSec{}s. Scheduling overhead grows with $k$ through $k{=}8$ and then plateaus, because each dispatch fans out to more predictors and is bounded by the slowest reply.

Randomized probing also mitigates herding. When many bursty requests arrive within one round trip, a deterministic Fanout scheduler can compute identical predictions and overload the same instance, while a power-of-$k$ policy breaks this correlation by randomizing each request's candidate set. This helps the sampled variants maintain lower E2E latency and higher throughput than Fanout in Figure~\ref{fig:po2_comparison}. In this sweep, $k{=}4$ gives the best empirical tradeoff: it has the highest capacity and throughput, while retaining the second-lowest E2E latency and scheduling overhead. We nevertheless use Po2 as the default in the main evaluation to follow the standard power-of-two-choices convention without tuning $k$ on the evaluation workload.

% Oracle-length upper bounds (§6.7): est vs oracle × Po2 vs Fanout
\subsection{Oracle-Length Upper Bounds}\label{section:oracle_bounds}

Astrolabe uses a RoBERTa length estimator (MAPE 24.4\%, \S\ref{section:predicting accuracy}). Replacing it with ground-truth (``oracle'') lengths gives an upper bound on the gains Po2 and Fanout can obtain from a perfect length predictor. Length-controllable workloads do exist, such as constrained decoding~\cite{geng2025jsonschemabench,beurerkellner2022lmql,willard2023efficient} and length-adaptive applications~\cite{reddy2025bellman}, which can provide the exact length of each request at dispatch time. In addition, applications with standard output formats can have highly deterministic output lengths even without model-based prediction, making the oracle mode realistic for these scenarios.
\begin{figure}[!tb]
    \centering
    \includegraphics[width=0.90\linewidth]{figure/revision/oracle_comparison.png}
    \caption{Oracle-length upper bounds}
    \Description{Capacity, mean E2E latency, mean TTFT, and scheduling overhead comparing estimated and oracle lengths for Po2 and Fanout.}
    \label{fig:oracle_comparison}
\end{figure}

Figure~\ref{fig:oracle_comparison} compares the four variants across capacity, mean E2E latency, mean TTFT, and scheduling overhead. Oracle adds $\sim$0.8\,QPS for Po2 (\capPoTwoEst{}\,$\to$\,\capPoTwoOracle{}) and $\sim$0.9\,QPS for Fanout (\capFanoutEst{}\,$\to$\,\capFanoutOracle{}) at SLO\,=\,\sloSec{}s, a $\sim$3\% capacity gain. The default Astrolabe Po2 variant has the lowest scheduling overhead, while Fanout-oracle has the highest. Based on our observation, one candidate explanation is that estimated lengths are more uniform than actual lengths, which may increase the simulator's cache hit rate for repeated batch-latency lookups; the oracle variants expose more exact length diversity and therefore incur higher simulation overhead. Even for Fanout-oracle, however, scheduling overhead stays below 3\% of total E2E latency in this sweep and tends to decrease once capacity is reached, as E2E latency then spikes more rapidly.

% Burstiness - Po2 vs Llumnix across arrival variance
\subsection{Ablation Study Under Bursty Workloads}\label{section:burstiness}
Real-world LLM traffic exhibits bursty patterns~\cite{burstgpt}. We vary burstiness with gamma-distributed arrivals and set the shape parameter $\alpha$ to 0.25, 0.5, 1.0, and 2.0. Lower $\alpha$ produces burstier, higher-variance inter-arrival times; at $\alpha{=}1.0$, the arrival process reduces to the Poisson case with exponential inter-arrival times. All experiments run at QPS=32 on the A30 cluster against Llumnix-, and we record E2E latencies.

\begin{figure}[!tb]
    \centering
    \includegraphics[width=0.90\linewidth]{figure/revision/burstiness_lines.png}
    \caption{Latency under bursty workloads}
    \Description{E2E latency under gamma arrival processes with varying burstiness for Astrolabe and Llumnix-.}
    \label{fig:burstiness}
\end{figure}

Figure~\ref{fig:burstiness} shows Astrolabe matching or beating Llumnix- across all four levels. Under high burstiness ($\alpha{=}0.25$), Astrolabe's mean E2E latency is about 9\% lower than Llumnix-; under near-regular arrivals ($\alpha{=}2.0$), the gap narrows but Astrolabe retains a small edge. This pattern is consistent with randomized probing decorrelating placement decisions under bursts, reducing request herding while keeping per-request probing cost fixed at two predictors. Moving from $\alpha{=}2.0$ to $\alpha{=}0.25$ raises mean E2E by $\sim$13\% for Astrolabe vs $\sim$22\% for Llumnix-. Thus, Astrolabe's Po2 dispatch is more robust to burstiness because randomized probing decorrelates placement decisions.

% Prediction-error sensitivity heatmap (Po2-est × length_err × latency_err)
\subsection{Sensitivity to Prediction Errors}\label{section:error_sensitivity}

A key design question is whether length-prediction error (24.4\%, Table~\ref{table:length_prediction}) and simulator error (10--15\%, Figure~\ref{fig:latency_prediction}) undermine dispatch quality. We answer this by injecting noise well beyond both: $v' = v \cdot (1 + \epsilon)$ with $\epsilon \sim \mathcal{N}(0, \sigma)$, clamped to $[0.1v, 5.0v]$ to avoid invalid values (e.g., negative lengths). We inject noise independently into the length estimate and each instance's simulated latency. We sweep $\sigma$ over 0\%, 25\%, 50\%, and 100\% on each axis separately, producing a 4$\times$4 grid at QPS=32 with 10k requests per cell to demonstrate the robustness of Po2 dispatch across a wide range of error profiles, including extreme cases where noise dominates the signal.

\begin{figure}[!tb]
    \centering
    \includegraphics[width=0.90\linewidth]{figure/revision/prediction_error_heatmap.png}
    \caption{Sensitivity of Latency to Prediction Error}
    \Description{Heatmaps of mean and P99 E2E degradation as length-estimation and latency-prediction noise increase.}
    \label{fig:error_heatmap}
\end{figure}

Figure~\ref{fig:error_heatmap} shows mean (left) and P99 (right) E2E degradation vs the noise-free baseline. The largest mean degradation is +6.3\% at (100\%,\,25\%), while the extreme (100\%,\,100\%) cell is +5.1\%; most cells remain within the $\pm$4\% single-deploy noise floor. Figure~\ref{fig:error_heatmap} overlays Llumnix-'s degradation (+4.0\% mean, +4.1\% P99) as a colorbar reference. Astrolabe stays below this line at moderate noise; at high length-error settings, several mean-latency cells exceed it modestly, but P99 degradation remains below Llumnix- in every cell (maximum +3.7\%). Thus predictor-guided dispatch retains value across the injected-noise range. Three factors keep Po2 robust: (i) dispatch uses \emph{relative} ranking between the two sampled probes, so common-mode length error cancels and only large per-instance latency noise can flip a ranking; (ii) underlying latency gaps between loaded and idle instances are usually 2--10\,s, so moderate noise rarely reverses orderings; (iii) Po2's random probe set prevents one bad ranking from systematically steering later requests to the same overloaded instance, which is the same randomization principle that mitigates herding in~\S\ref{section:po2_ablation}. Length underestimation is also mitigated by the auto-extension (\S\ref{section:model_instance}), which applies the same 10-token extension rule across candidate simulations.

% CPU overhead — Po2 vs Fanout, direct measurement
\subsection{CPU Overhead: Po2 vs Fanout}\label{section:cpu_overhead}

Astrolabe's Predictor runs CPU-side simulation off the GPU-bound inference path. The control-plane cost is governed by the number of predictor probes per dispatch: Fanout queries all 12 predictors, while Astrolabe's default Po2 policy queries only two. We measure per-predictor CPU utilization and RSS memory on the A30 cluster.
\begin{figure}[!tb]
    \centering
    \includegraphics[width=0.90\linewidth]{figure/revision/cpu_overhead.png}
    \caption{Per-predictor CPU utilization and memory across QPS}
    \Description{Bar charts of mean CPU utilization and RSS memory per predictor across QPS for Po2 and Fanout.}
    \label{fig:cpu_overhead}
\end{figure}
Figure~\ref{fig:cpu_overhead} shows that at QPS=20, mean per-predictor CPU is \cpuPoTwoMeanAtTwenty{}\% for Po2 and \cpuFanoutMeanAtTwenty{}\% for Fanout, a \cpuPoTwoVsFanoutRatio{}$\times$ reduction. The gap persists across QPS 20, 24, 28, 32, and 36 because Fanout sends every request to all 12 predictors, whereas Po2 sends each request to only two. In contrast, per-predictor RSS is stable across QPS within each variant: at QPS=20, Po2 uses 1225\,MB and Fanout uses 1773\,MB. Taken together, these measurements show that the scaling difference is CPU-bound: predictor RSS is a stable fixed cost, while Fanout's aggregate predictor CPU grows with cluster size because it probes every instance. Po2 keeps the per-request control-plane fanout constant at two probes.


\subsection{Generality Study}\label{section:general_study}
\begin{table}[htb]
\caption{Capacity across workload/configuration changes. bs denotes batch size and cs denotes chunk size.}
\setlength{\tabcolsep}{4pt}
\resizebox{\linewidth}{!}{%
\begin{tabular}{lcccc}
\toprule
Scheduler    & bs=24 & cs=2048 & Qwen2-7B & BurstGPT \\
\midrule
Astrolabe-oracle  & 29.1 & 31.9 & \capQwenPoTwoOracle{} & \capBurstGPTPoTwoOracle{} \\
Astrolabe         & 28.5 & 31.6 & \capQwenPoTwoEst{}    & /    \\
Llumnix-          & 26.6 & 31.3 & \capQwenLlumnix{}     & 61.6 \\
Gain (est/oracle) & 7.1\%/9.4\% & 1.0\%/1.9\% & 6.0\%/12.6\% & /\,4.7\% \\
\bottomrule
\end{tabular}%
}
\label{tab:scheduler_perf_scaled}
\end{table}

Dynamic changes in LLM inference settings (such as batch size, chunk size, model, or dataset) are typically ignored by heuristic schedulers' static rules and cause performance drift (\S\ref{section:scheduling_issues}). In contrast, Astrolabe's predictor simulates the specific backend configuration, so changes are absorbed automatically. Table~\ref{tab:scheduler_perf_scaled} compares Astrolabe against Llumnix- across four variations to validate this hypothesis.

Across these variations: (1)~\emph{batch/chunk}---with batch size reduced 48$\to$24, Astrolabe reaches 28.5 QPS vs Llumnix- 26.6 QPS (+7.1\%) and Astrolabe-oracle reaches 29.1 QPS (+9.4\%); with chunk size raised 512$\to$2048, the gap narrows to 1.0\%/1.9\%, which is similar to the main-experiment capacity difference, because the larger chunked-prefill step changes the predictor's inputs less disruptively than a batch-size change. (2)~\emph{model: Qwen2-7B}---after retraining the RoBERTa length model on ShareGPT tokenized with Qwen, Astrolabe reaches \capQwenPoTwoEst{}~QPS vs Llumnix- \capQwenLlumnix{}~QPS (+6.0\%, +12.6\% with oracle lengths). (3)~\emph{dataset: BurstGPT}~\cite{burstgpt}---the trace provides lengths without prompt content, which prevents retraining the length predictor,
so we compare only oracle variants: Astrolabe-oracle reaches \capBurstGPTPoTwoOracle{}~QPS vs Llumnix- 61.6~QPS (+4.7\%). Across all four, Astrolabe or its oracle variant matches or outperforms Llumnix-, confirming that prediction-guided randomized dispatch is robust to backend, model, and workload changes.


% \subsection{Auto Provisioning} \label{section:auto_provision}
% As discussed in \S\ref{section:scheduling_issues}, provisioning based on runtime metrics could suffer from asynchronous cold start. Furthermore, accurate predicted metrics can be utilized for provisioning in a preemptive manner. To demonstrate this, we implement two simple latency-based auto-provisioning strategies. The \emph{preempt} strategy adds an instance when predicted latency (using ground-truth length with 15\% prediction accuracy in \S \ref{section:predicting accuracy}) reaches 70 seconds, while \emph{relief} strategy means provisioning only when actual latency hits the same threshold. We conduct experiments starting with six available instances under a QPS of 24. We also establish a baseline with a sufficient static cluster of 10 instances to allow all requests completed in time.

%  \begin{figure*}[!tb]
%         \centering % <-- added
%         \includegraphics[width=\linewidth]{figure/exp_plots/cluster_metrics/auto_provision.png}
%          \caption{Auto Provisioning with Different Strategy Under QPS=24}
%          \label{fig:auto_provsion}
%          \subcaption*{Dotted lines at left-bottom figure indicate the change of scheduling instances when triggering provisioning.}
% \end{figure*}

% Figure ~\ref{fig:auto_provsion} presents the results, with the first column displaying the dot plot of request latencies and the changes of number of instances in the clusters along with Query IDs sent in order. The second column illustrates changes in the GPU free memory blocks of the clusters and the variances across instances

% These results support our hypothesis that a \emph{relief} strategy causes over-provisioning.  As newly added hosts are unable to alleviate queued requests, which could trigger further provisioning, all backup instances become quickly exhausted. This situation causes memory imbalances and GPU resource wastage as more instances get provisioned in rush. In contrast, the \emph{preempt} strategy activates earlier, resulting in smoother changes in cluster size and ultimately provisioning only the necessary instances. Consequently, the second column shows lower variance and higher memory utilization over time. Also, the \emph{preempt} strategy ultimately leads to impressive results in the latency dot plots, reducing P99 latency by 20.1\% (from 89.2 to 71.2) and decreasing the number of requests exceeding 70 seconds by 81\% (from 627 to 129), even with fewer instances (from 12 to 10). We believe that the improvements from predictive-based provisioning can persist when applied to or extended into other complex and practical provisioning strategies in large-scale, real-world cluster management.  


% \section{Future Work} 
% Our current implementation already supports key components similarly to other LLM scheduling frameworks. As part of our future work, we plan to add features to align with the new key trends in LLM serving. While P-D disaggregation has not yet been integrated, we consider its support should be feasible. Similar to LLM-d~\cite{githubGitHubLlmdllmd}, it would involve dedicated schedulers for prefill and decode phases with inter-phase cache transferring enabled. Given that inter-instance cache transfer remains under active development within inference frameworks, and that aggregated architectures continue to achieve competitive performance through optimizations such as chunked-prefill \cite{agrawal2023sarathiefficientllminference}, optimized kernels \cite{podAttention}, and GPU multiplexing \cite{gao2025duetserveharmonizingprefilldecode}, we consider a comprehensive implementation of P-D disaggregation to be outside the scope of current work and reserve this exploration for future research.

% Additionally, cache transfer enables the exploration of Block Scheduler combined with dynamic rebalancing. Nevertheless, we anticipate that when integrating the above features Block's performance advantages will persist, as the fundamental scheduling challenges discussed in \S\ref{section:scheduling_issues} remain pertinent in disaggregated settings. Furthermore, incorporating prefix caching~\cite{vllmAutomaticPrefix} can enhance simulator accuracy for real-world applications involving multi-turn conversations and repetitive prompts~\cite{sglang}. Since it has a negligible impact on experimental setup with de-duplicated dataset, we leave it for further exploration.

% In addition to integrating the newly supported internal features from vLLM into Block, we plan to evaluate Block against larger models using other multi-GPU equipped hosts and model parallelism. As discussed in Section \ref{section:experiment_setting}, Vidur's evaluation indicates that larger models within the same classes achieve higher simulation accuracies. This could result in greater gains and more significant positive impacts from Block. We will also schedule future updates with the revised testbed setups in the future.

\section{Conclusion}
Astrolabe combines response-length estimation, simulation-based latency prediction, and randomized dispatch for multi-instance LLM serving. Across our experiments, prediction-guided one-shot dispatch matches or outperforms migration-based rebalancing while avoiding KV-cache transfer overhead.

\bibliographystyle{ACM-Reference-Format}
\bibliography{reference}


\end{document}
